
\label{sec:result}
\subsection{Code Coverage of MIST (RQ1)}

\textbf{Experimental Design.} To answer this research question, we evaluate MIST on three DBMSs (DuckDB, PostgreSQL, SQLite) using four LLMs (Qwen2.5-7B, Llama3.1-8B, Qwen2.5-14B, Qwen2.5-32B).

\textbf{Results Analysis.} Table~\ref{table:rq1} shows the coverage comparison between the baseline and MIST. Based on the results, we make the following observations.

\textbf{(1) MIST consistently improves coverage across all DBMSs and coverage metrics.} Across the three DBMSs, MIST achieves substantial coverage improvements over the baseline. For DuckDB, MIST improves line coverage by an average of 49.0\%, function coverage by 33.6\%, and branch coverage by 45.8\%. For PostgreSQL, the improvements are 34.0\% in line coverage, 35.1\% in function coverage, and 35.6\% in branch coverage. For SQLite, MIST achieves improvements of 47.0\% in line coverage, 28.3\% in function coverage, and 57.7\% in branch coverage. These consistent improvements across diverse DBMS architectures, ranging from in-process analytical (DuckDB) to client-server relational (PostgreSQL) and embedded transactional (SQLite), demonstrate that MIST is effective regardless of the underlying database design.

\textbf{(2) MIST shows strong generalization across different LLM sizes and families.} The effectiveness of MIST is consistent across all four evaluated LLMs, regardless of parameter scale or model family. For the Qwen family, MIST achieves branch coverage improvements ranging from 31.2\% to 69.2\% on DuckDB, 31.5\% to 54.2\% on PostgreSQL, and 39.1\% to 92.9\% on SQLite. For Llama3.1-8B, MIST achieves improvements of 38.8\%, 21.3\%, and 52.1\% on the three DBMSs, respectively. Notably, smaller models with MIST often outperform larger baseline models. For example, Qwen2.5-7B with MIST achieves 42.0\% line coverage on DuckDB, surpassing the baseline of Qwen2.5-32B (25.8\%). This demonstrates that MIST is more critical than raw model capacity, making it practical for resource-constrained industrial environments where deploying large models may not be feasible.

% \begin{tcolorbox}[breakable,width=\linewidth-2pt,boxrule=0pt,top=2pt, bottom=2pt, left=4pt,right=4pt, colback=gray!15,colframe=gray!15]
% \textbf{Answer to RQ1:} MIST consistently improves code coverage across all three DBMSs and four LLMs, achieving average improvements of 43.3\% in line coverage, 32.3\% in function coverage, and 46.4\% in branch coverage compared to the baseline approach. The improvements demonstrate strong generalization across different DBMS architectures and LLM sizes.
% \end{tcolorbox}

\noindent
\fcolorbox{gray!15}{gray!15}{
\parbox{\dimexpr\linewidth-2\fboxsep-2\fboxrule}{
\textbf{Answer to RQ1:} MIST consistently improves code coverage across all three DBMSs and four LLMs, achieving average improvements of 43.3\% in line coverage, 32.3\% in function coverage, and 46.4\% in branch coverage compared to the baseline approach. The improvements demonstrate strong generalization across different DBMS architectures and LLM sizes.
}}


% Please add the following required packages to your document preamble:
% \usepackage{multirow}
% \usepackage[table,xcdraw]{xcolor}
% Beamer presentation requires \usepackage{colortbl} instead of \usepackage[table,xcdraw]{xcolor}
% \begin{table*}[t]
% \caption{xxx}
% \label{table:rq2}
% \begin{tabular}{c|c|cccc|cccc|cccc}
% \toprule
%  &  & \multicolumn{4}{c|}{DuckDB} & \multicolumn{4}{c|}{\textbf{PostgreSQL}} & \multicolumn{4}{c}{\textbf{SQlite}} \\ \cline{3-14} 
% \multirow{-2}{*}{\textbf{Model}} & \multirow{-2}{*}{\textbf{Approach}} & \textbf{Parser} & \textbf{Optimizer} & \textbf{Executor} & \textbf{Storage} & \textbf{Parser} & \textbf{Optimizer} & \textbf{Executor} & \textbf{Storage} & \textbf{Parser} & \textbf{Optimizer} & \textbf{Executor} & \textbf{Storage} \\ \midrule
%  & Base &  &  &  &  &  &  &  &  &  &  &  &  \\
% \multirow{-2}{*}{Qwen2.5-7B} & \cellcolor[HTML]{ECF4FF}\textbf{MIST} & \cellcolor[HTML]{ECF4FF}\textbf{} & \cellcolor[HTML]{ECF4FF}\textbf{} & \cellcolor[HTML]{ECF4FF}\textbf{} & \cellcolor[HTML]{ECF4FF} & \cellcolor[HTML]{ECF4FF}\textbf{} & \cellcolor[HTML]{ECF4FF}\textbf{} & \cellcolor[HTML]{ECF4FF}\textbf{} & \cellcolor[HTML]{ECF4FF} & \cellcolor[HTML]{ECF4FF}\textbf{} & \cellcolor[HTML]{ECF4FF}\textbf{} & \cellcolor[HTML]{ECF4FF}\textbf{} & \cellcolor[HTML]{ECF4FF} \\ \midrule
%  & Base &  &  &  &  &  &  &  &  &  &  &  &  \\
% \multirow{-2}{*}{Llama3.1-8B} & \cellcolor[HTML]{ECF4FF}\textbf{MIST} & \cellcolor[HTML]{ECF4FF}\textbf{} & \cellcolor[HTML]{ECF4FF}\textbf{} & \cellcolor[HTML]{ECF4FF}\textbf{} & \cellcolor[HTML]{ECF4FF} & \cellcolor[HTML]{ECF4FF}\textbf{} & \cellcolor[HTML]{ECF4FF}\textbf{} & \cellcolor[HTML]{ECF4FF}\textbf{} & \cellcolor[HTML]{ECF4FF} & \cellcolor[HTML]{ECF4FF}\textbf{} & \cellcolor[HTML]{ECF4FF}\textbf{} & \cellcolor[HTML]{ECF4FF}\textbf{} & \cellcolor[HTML]{ECF4FF} \\ \midrule
%  & Base &  &  &  &  &  &  &  &  &  &  &  &  \\
% \multirow{-2}{*}{Qwen2.5-14B} & \cellcolor[HTML]{ECF4FF}\textbf{MIST} & \cellcolor[HTML]{ECF4FF}\textbf{} & \cellcolor[HTML]{ECF4FF}\textbf{} & \cellcolor[HTML]{ECF4FF}\textbf{} & \cellcolor[HTML]{ECF4FF} & \cellcolor[HTML]{ECF4FF}\textbf{} & \cellcolor[HTML]{ECF4FF}\textbf{} & \cellcolor[HTML]{ECF4FF}\textbf{} & \cellcolor[HTML]{ECF4FF} & \cellcolor[HTML]{ECF4FF}\textbf{} & \cellcolor[HTML]{ECF4FF}\textbf{} & \cellcolor[HTML]{ECF4FF}\textbf{} & \cellcolor[HTML]{ECF4FF} \\ \midrule
%  & Base &  &  &  &  &  &  &  &  &  &  &  &  \\
% \multirow{-2}{*}{Qwen2.5-14B} & \cellcolor[HTML]{ECF4FF}\textbf{MIST} & \cellcolor[HTML]{ECF4FF}\textbf{} & \cellcolor[HTML]{ECF4FF}\textbf{} & \cellcolor[HTML]{ECF4FF}\textbf{} & \cellcolor[HTML]{ECF4FF} & \cellcolor[HTML]{ECF4FF}\textbf{} & \cellcolor[HTML]{ECF4FF}\textbf{} & \cellcolor[HTML]{ECF4FF}\textbf{} & \cellcolor[HTML]{ECF4FF} & \cellcolor[HTML]{ECF4FF}\textbf{} & \cellcolor[HTML]{ECF4FF}\textbf{} & \cellcolor[HTML]{ECF4FF}\textbf{} & \cellcolor[HTML]{ECF4FF} \\ \bottomrule
% \end{tabular}
% \end{table*}

% Please add the following required packages to your document preamble:
% \usepackage[table,xcdraw]{xcolor}
% Beamer presentation requires \usepackage{colortbl} instead of \usepackage[table,xcdraw]{xcolor}
% \begin{table}[t]
% \caption{Module-level line coverage comparison between Fuzz4All  and MIST.}
% \label{table:rq2}
% \begin{tabular}{ccccc}
% \toprule
% \multicolumn{1}{c|}{\textbf{Approach}} & \textbf{Parser} & \textbf{Optimizer} & \textbf{Executor} & \textbf{Storage} \\ \midrule
% \multicolumn{5}{c}{\cellcolor[HTML]{ECF4FF}\textbf{DuckDB}} \\ \midrule
% \multicolumn{1}{c|}{Fuzz4All (Qwen2.5-7B)} & 26.6\% & 51.6\% & 31.0\% & 20.0\% \\
% \multicolumn{1}{c|}{\textit{\textbf{+ MIST}}} & \textbf{43.2\%} & \textbf{68.4\%} & \textbf{42.2\%} & \textbf{31.3\%} \\ \midrule
% \multicolumn{1}{c|}{Llama3.1-8B} & 30.1\% & 52.8\% & 35.1\% & 22.7\% \\
% \multicolumn{1}{c|}{\textbf{+ MIST}} & \textbf{48.9\%} & \textbf{69.3\%} & \textbf{47.3\%} & \textbf{32.9\%} \\ \midrule
% \multicolumn{1}{c|}{Qwen2.5-14B} & 24.2\% & 43.4\% & 22.9\% & 22.1\%  \\
% \multicolumn{1}{c|}{\textbf{+ MIST}} & \textbf{45.1\%} & \textbf{67.0\%} & \textbf{46.6\%} & \textbf{31.8\%} \\ \midrule
% \multicolumn{1}{c|}{Qwen2.5-32B}& 23.7\% & 45.6\% &  26.4\% & 22.6\%  \\
% \multicolumn{1}{c|}{\textbf{+ MIST}} & \textbf{43.0\%} & \textbf{66.9\%} & \textbf{45.2\%} & \textbf{31.8\%} \\ \midrule
% \multicolumn{5}{c}{\cellcolor[HTML]{ECF4FF}\textbf{PostgreSQL}} \\ \midrule
% \multicolumn{1}{c|}{Qwen2.5-7B} & 36.1\% & 50.8\% & 37.5\% & 31.8\% \\
% \multicolumn{1}{c|}{\textit{\textbf{+ MIST}}} & \textbf{41.0\%} & \textbf{59.0\%} & \textbf{43.1\%} & \textbf{36.4\%} \\ \midrule
% \multicolumn{1}{c|}{Llama3.1-8B} & 32.9\% & 47.3\% & 33.5\% &  30.9\% \\
% \multicolumn{1}{c|}{\textit{\textbf{+ MIST}}} & \textbf{39.9\%} & \textbf{57.4\%} & \textbf{39.1\%} & \textbf{36.3\%} \\ \midrule
% \multicolumn{1}{c|}{Qwen2.5-14B} & 32.2\% & 44.6\%  & 30.6\% &  30.5\%\\
% \multicolumn{1}{c|}{\textit{\textbf{+ MIST}}} & \textbf{41.5\%} & \textbf{62.6\%} & \textbf{45.8\%} & \textbf{38.1\%} \\ \midrule
% \multicolumn{1}{c|}{Qwen2.5-32B} & 29.6\% & 38.4\% & 25.3\% & 27.9\%  \\
% \multicolumn{1}{c|}{\textit{\textbf{+ MIST}}} & \textbf{41.2\%} & \textbf{63.4\%} & \textbf{45.9\%} & \textbf{38.1\%} \\ \bottomrule
% % 注意这里caption里要写一下为啥没有SQLite的模块覆盖率
% \end{tabular}
% \end{table}

\begin{table}[t]
\caption{Module-level line coverage comparison between baseline and MIST.}
\label{table:rq2}
\small
\centering
\begin{tabular}{c|c|cccc}
\toprule
\textbf{Model} & \textbf{Approach} & \textbf{Parser} & \textbf{Optimizer} & \textbf{Executor} & \textbf{Storage} \\ 
\midrule

\multicolumn{6}{c}{\cellcolor[HTML]{ECF4FF}\textbf{DuckDB}} \\ \midrule

\multirow{2}{*}{\begin{tabular}[c]{@{}c@{}}Qwen2.5\\ 7B\end{tabular}}
  & Fuzz4All & 26.6\% & 51.6\% & 31.0\% & 20.0\% \\
  & \textbf{MIST} & \textbf{43.2\%} & \textbf{68.4\%} & \textbf{42.2\%} & \textbf{31.3\%} \\ \midrule

\multirow{2}{*}{\begin{tabular}[c]{@{}c@{}}Llama3.1\\ 8B\end{tabular}}
  & Fuzz4All  & 30.1\% & 52.8\% & 35.1\% & 22.7\% \\
  & \textbf{MIST} & \textbf{48.9\%} & \textbf{69.3\%} & \textbf{47.3\%} & \textbf{32.9\%} \\ \midrule

\multirow{2}{*}{\begin{tabular}[c]{@{}c@{}}Qwen2.5\\ 14B\end{tabular}}
  & Fuzz4All  & 24.2\% & 43.4\% & 22.9\% & 22.1\% \\
  & \textbf{MIST} & \textbf{45.1\%} & \textbf{67.0\%} & \textbf{46.6\%} & \textbf{31.8\%} \\ \midrule

\multirow{2}{*}{\begin{tabular}[c]{@{}c@{}}Qwen2.5\\ 32B\end{tabular}}
  & Fuzz4All  & 23.7\% & 45.6\% & 26.4\% & 22.6\% \\
  & \textbf{MIST} & \textbf{43.0\%} & \textbf{66.9\%} & \textbf{45.2\%} & \textbf{31.8\%} \\ \midrule

\multicolumn{6}{c}{\cellcolor[HTML]{ECF4FF}\textbf{PostgreSQL}} \\ \midrule

\multirow{2}{*}{\begin{tabular}[c]{@{}c@{}}Qwen2.5\\ 7B\end{tabular}}
  & Fuzz4All  & 36.1\% & 50.8\% & 37.5\% & 31.8\% \\
  & \textbf{MIST} & \textbf{41.0\%} & \textbf{59.0\%} & \textbf{43.1\%} & \textbf{36.4\%} \\ \midrule

\multirow{2}{*}{\begin{tabular}[c]{@{}c@{}}Llama3.1\\ 8B\end{tabular}}
  & Fuzz4All  & 32.9\% & 47.3\% & 33.5\% &  30.9\% \\
  & \textbf{MIST} & \textbf{39.9\%} & \textbf{57.4\%} & \textbf{39.1\%} & \textbf{36.3\%} \\ \midrule

\multirow{2}{*}{\begin{tabular}[c]{@{}c@{}}Qwen2.5\\ 14B\end{tabular}} 
  & Fuzz4All  & 32.2\% & 44.6\%  & 30.6\% &  30.5\% \\
  & \textbf{MIST} & \textbf{41.5\%} & \textbf{62.6\%} & \textbf{45.8\%} & \textbf{38.1\%} \\ \midrule

\multirow{2}{*}{\begin{tabular}[c]{@{}c@{}}Qwen2.5\\ 32B\end{tabular}}
  & Fuzz4All  & 29.6\% & 38.4\% & 25.3\% & 27.9\% \\
  & \textbf{MIST} & \textbf{41.2\%} & \textbf{63.4\%} & \textbf{45.9\%} & \textbf{38.1\%} \\ 

\bottomrule
\end{tabular}
\end{table}

% Please add the following required packages to your document preamble:
% \usepackage{multirow}
\begin{table*}[t]
\caption{Ablation study results showing the impact of different components in MIST. Each row represents a different combination of synthetization strategy and mutation strategy.}
\label{table:rq3}
\begin{tabular}{c|c|ccc|ccc|ccc}
\toprule
\multirow{2}{*}{\textbf{Synthetization}} & \multirow{2}{*}{\textbf{Mutation}} & \multicolumn{3}{c|}{\textbf{DuckDB}} & \multicolumn{3}{c|}{\textbf{PostgreSQL}} & \multicolumn{3}{c}{\textbf{SQLite}} \\ \cline{3-11} 
 &  & \textbf{Line} & \textbf{Function} & \textbf{Branch} & \textbf{Line} & \textbf{Function} & \textbf{Branch} & \textbf{Line} & \textbf{Function} & \textbf{Branch} \\ \midrule
\begin{tabular}[c]{@{}c@{}}Simple \\ Instruction\end{tabular} & \XSolidBrush  & 29.9\% & 29.7\%  & 17.0\% & 23.1\% & 25.2\% & 16.2\% & 29.9\% & 38.9\%  & 22.5\%  \\ \midrule
\begin{tabular}[c]{@{}c@{}}Random\\ Feature\end{tabular} & \XSolidBrush & 32.3\% & 31.7\% & 18.7\% & 25.2\% & 30.1\% & 16.9\% & 32.8\% & 41.7\% & 24.6\% \\ \midrule
\begin{tabular}[c]{@{}c@{}}Hierarchical\\ Feature \\ (MIST)\end{tabular} & \XSolidBrush & 36.2\% & 34.7\% & 20.6\% & 27.4\% & 32.1\% & 17.8\% & 35.5\% & 43.2\% & 26.3\%
 \\ \midrule \midrule
\begin{tabular}[c]{@{}c@{}}Hierarchical\\ Feature \\ (MIST)\end{tabular} & \begin{tabular}[c]{@{}c@{}}Simple \\ Instruction\end{tabular} & 37.8\% & 35.7\% & 21.6\% & 28.3\% & 33.9\% & 18.1\%  & 36.9\% & 44.2\% & 27.5\% \\ \midrule
\begin{tabular}[c]{@{}c@{}}Hierarchical\\ Feature \\ (MIST)\end{tabular} & \begin{tabular}[c]{@{}c@{}}Random\\ Rule\end{tabular} & 38.7\% & 36.7\% & 22.4\% & 28.8\% & 35.4\% & 19.5\% & 37.5\% & 44.8\% & 29.2\% \\ \midrule
\textbf{\begin{tabular}[c]{@{}c@{}}Hierarchical\\ Feature \\ (MIST)\end{tabular}} & \textbf{\begin{tabular}[c]{@{}c@{}}MCTS\\ Rule \\ (MIST)\end{tabular}} & \textbf{42.0\%} & \textbf{38.2\%} & \textbf{24.5\%} & \textbf{30.3\%} & \textbf{38.6\%} & \textbf{21.3\%} & \textbf{39.7\%} & \textbf{46.1\%} & \textbf{31.3\%} \\ 

\bottomrule
\end{tabular}
\end{table*}

\subsection{Module-Level Coverage of MIST (RQ2)}
\noindent \textbf{Experimental Design.} To answer this research question, we analyze the line coverage distribution across four critical DBMS modules: \textit{Parser}, \textit{Optimizer}, \textit{Executor}, and \textit{Storage}. We focus on DuckDB and PostgreSQL, as SQLite's monolithic architecture does not provide clear module boundaries for fine-grained analysis.

\noindent \textbf{Results Analysis.} Table~\ref{table:rq2} shows the module-level line coverage comparison between the baseline and MIST. Based on the results, we make the following observations.

\textbf{(1) MIST achieves substantial improvements across all four modules.} For DuckDB, MIST improves \textit{Parser} coverage by an average of 73.2\%, \textit{Optimizer} coverage by 41.3\%, \textit{Executor} coverage by 61.4\%, and \textit{Storage} coverage by 46.5\% across the four LLMs. For PostgreSQL, the improvements are 25.8\% in \textit{Parser}, 35.8\% in \textit{Optimizer}, 40.7\% in \textit{Executor}, and 23.4\% in \textit{Storage}. These comprehensive improvements demonstrate that MIST effectively exercises execution paths across the entire DBMS pipeline, from front-end parsing to back-end storage operations. Notably, even the \textit{Storage} module, which handles low-level data access and buffer management, shows consistent gains of 46.5\% and 23.4\% for DuckDB and PostgreSQL, respectively, indicating that MIST's test cases successfully trigger diverse storage access patterns.

\textbf{(2) The Optimizer module achieves the highest absolute coverage across both DBMSs.} Among the four modules, the \textit{Optimizer} consistently achieves the highest absolute coverage values with MIST. In DuckDB, MIST reaches an average of 67.9\% \textit{Optimizer} coverage, compared to 48.4\% in the baseline. In PostgreSQL, MIST achieves 60.6\% \textit{Optimizer} coverage, up from 45.3\% in the baseline. This is particularly because the \textit{Optimizer} is one of the most complex components in modern DBMSs, responsible for query plan generation, cost estimation, and selecting optimal execution strategies. The high coverage indicates that MIST's feature-guided generation and MCTS-based mutation effectively trigger diverse optimization logic. For example, Llama3.1-8B with MIST achieves 69.3\% \textit{Optimizer} coverage on DuckDB, demonstrating that even smaller models can thoroughly exercise complex optimization paths when guided by MIST.

\textbf{(3) Different modules benefit differently from MIST, reflecting their architectural characteristics.} The magnitude of improvement varies across modules, with \textit{Parser} and \textit{Executor} showing larger relative gains in DuckDB (73.2\% and 61.4\%), while PostgreSQL shows more balanced improvements across all modules (23.4\% to 40.7\%). This difference reflects DuckDB's specialized design for analytical workloads, where \textit{Parser} and \textit{Executor} handle complex operations like window functions, CTEs, and vectorized execution. In contrast, PostgreSQL's mature and extensively tested codebase shows more uniform improvements, indicating that MIST discovers previously untested paths across all components rather than concentrating on specific modules.

% \begin{tcolorbox}[breakable,width=\linewidth-2pt,boxrule=0pt,top=2pt, bottom=2pt, left=4pt,right=4pt, colback=gray!15,colframe=gray!15]
% \textbf{Answer to RQ2:} MIST achieves substantial improvements across all DBMS modules, with the \textit{Optimizer} module showing the highest absolute coverage (average 67.9\% for DuckDB and 60.6\% for PostgreSQL).
% \end{tcolorbox}

\noindent
\fcolorbox{gray!15}{gray!15}{
\parbox{\dimexpr\linewidth-2\fboxsep-2\fboxrule}{
\textbf{Answer to RQ2:} MIST achieves substantial improvements across all DBMS modules, with the \textit{Optimizer} module showing the highest absolute coverage (average 67.9\% for DuckDB and 60.6\% for PostgreSQL).
}}



\subsection{Ablation Study (RQ3)}

\noindent \textbf{Experimental Design.} To answer this research question, we perform ablation studies by evaluating different combinations of synthetization and mutation strategies using Qwen2.5-7B. For synthetization, we compare three approaches:

\begin{itemize}
\item \textbf{Simple Instruction:} Uses basic prompts without feature guidance.
\item \textbf{Random Feature:} Randomly selects features without hierarchical structure.
\item \textbf{Hierarchical Feature:} Uses our hierarchical feature tree with diversity-aware sampling.
\end{itemize}

For mutation, we compare four approaches:

\begin{itemize}
\item \textbf{No Mutation:} Skips the mutation stage.
\item \textbf{Simple Instruction:} Uses basic prompts for mutation.
\item \textbf{Random Rule:} Randomly applies mutation rules.
\item \textbf{MCTS Rule:} Uses our MCTS-based mutation engine.
\end{itemize}

\noindent \textbf{Results Analysis.} Table~\ref{table:rq3} shows the coverage results of different component combinations. Based on the results, we make the following observations.

\textbf{(1) Hierarchical feature selection is the most effective synthetization strategy.} Comparing the three synthetization strategies without mutation, \textit{Hierarchical Feature} achieves branch coverage of 20.6\% on DuckDB, 17.8\% on PostgreSQL, and 26.3\% on SQLite. This outperforms \textit{Random Feature} by 10.2\%, 5.3\%, and 6.9\%, respectively, and \textit{Simple Instruction} by 21.2\%, 9.9\%, and 16.9\%, respectively. The superior performance of \textit{Hierarchical Feature} can be attributed to its ability to systematically sample features from the hierarchical tree while maintaining semantic coherence, ensuring that generated test cases exercise diverse DBMS functionalities. 
% In contrast, \textit{Random Feature} lacks structural guidance, and \textit{Simple Instruction} fails to leverage DBMS-specific knowledge, resulting in lower coverage.

\begin{figure*}[t]
    \centering
    \includegraphics[scale=0.55]{figure/case_study.pdf}
    \caption{Case study on test case generation using Qwen2.5-7B on SQLite.}
    \label{fig:case_study}
\end{figure*}


\textbf{(2) MCTS-based mutation consistently outperforms alternative mutation strategies.} Among the three mutation approaches applied to \textit{Hierarchical Feature}, \textit{MCTS Rule} achieves branch coverage of 24.5\% on DuckDB, 21.3\% on PostgreSQL, and 31.3\% on SQLite. This represents improvements of 13.4\%, 17.7\%, and 13.8\% over \textit{Simple Instruction}, and 9.4\%, 9.2\%, and 7.2\% over \textit{Random Rule}, respectively. The consistent superiority of \textit{MCTS Rule} demonstrates that its tree search mechanism effectively balances exploration of new mutation paths and exploitation of high-coverage seeds. \textit{Random Rule} performs better than \textit{Simple Instruction} but still falls short of \textit{MCTS Rule}, indicating that random selection cannot match the systematic search strategy provided by MCTS.

\textbf{(3) Removing either component leads to performance degradation, among which hierarchical feature selection has the greatest impact.} The complete MIST framework (Hierarchical Feature and MCTS Rule) achieves the highest coverage across all DBMSs. Removing the mutation stage (Hierarchical Feature and No Mutation) reduces branch coverage by 18.9\% on DuckDB, 19.7\% on PostgreSQL, and 19.0\% on SQLite. Replacing \textit{Hierarchical Feature} with \textit{Simple Instruction} while keeping \textit{MCTS Rule} would reduce coverage by 44.1\%, 31.5\%, and 39.1\% on the three DBMSs, respectively. These results demonstrate that both components are essential, with hierarchical feature selection being the more critical component.

% \begin{tcolorbox}[breakable,width=\linewidth-2pt,boxrule=0pt,top=2pt, bottom=2pt, left=4pt,right=4pt, colback=gray!15,colframe=gray!15]
% \textbf{Answer to RQ3:} Hierarchical feature selection outperforms baselines by 5.3\%-21.2\%, while MCTS-based mutation further improves coverage by 9.2\%-17.7\% over alternative mutation strategies. Removing either component leads to substantial performance degradation, with hierarchical feature selection having the greater impact, confirming the necessity of the complete framework.
% \end{tcolorbox}

\noindent
\fcolorbox{gray!15}{gray!15}{
\parbox{\dimexpr\linewidth-2\fboxsep-2\fboxrule}{
\textbf{Answer to RQ3:} Hierarchical feature selection outperforms baselines by 5.3\%-21.2\%, while MCTS-based mutation further improves coverage by 9.2\%-17.7\% over alternative mutation strategies. Removing either component leads to substantial performance degradation, with hierarchical feature selection having the greater impact, confirming the necessity of the complete framework.
}}
